\subsection{Схема 3: Восьмёркообразные узлы}

Колышка перекручивается несколько раз, а затем пробивается

\graphicspath{{\currentpath}}


\begin{tabular}{
	>{\centering\arraybackslash}m{1cm}
	>{\centering\arraybackslash}m{8cm}
	>{\centering\arraybackslash}m{3cm}
}
1 & \includegraphics[scale=2]{images/stividor-0.eps} & простой узел \\
2 & \includegraphics[scale=2]{images/stividor-1.eps} & восьмёрка \\
3 & \includegraphics[scale=2]{images/stividor-2.eps} & девятка \\
4 & \includegraphics[scale=2]{images/stividor-3.eps} & стивидорный узел
\end{tabular}


\subsection{Узлы похожие на восьмёрку, но с другой структурой}

Если сделать колышку и перекрутить её ещё на оборот, получится узел, похожий на стивидорный, но с нечётным числом пересечений.

Чтобы получить симметричный ему узел, нужно завязать простой узел, затем на сделать на нём колышку и пробить её ходовым концом.

Подобный узел может содержать только нечётное число пересечений
